%%=============================================================================
%% Conclusie
%%=============================================================================

\chapter{Conclusie}%
\label{ch:conclusie}

% TODO: Trek een duidelijke conclusie, in de vorm van een antwoord op de
% onderzoeksvra(a)g(en). Wat was jouw bijdrage aan het onderzoeksdomein en
% hoe biedt dit meerwaarde aan het vakgebied/doelgroep? 
% Reflecteer kritisch over het resultaat. In Engelse teksten wordt deze sectie
% ``Discussion'' genoemd. Had je deze uitkomst verwacht? Zijn er zaken die nog
% niet duidelijk zijn?
% Heeft het onderzoek geleid tot nieuwe vragen die uitnodigen tot verder 
%onderzoek?

Deze bachelorproef onderzocht in welke mate een AI-gedreven agent kan worden ingezet om dependency-updates te analyseren en te beheren, zodat het kernteam van Springbok dit proces efficiënter en met minder handmatig werk kan uitvoeren. De resultaten tonen aan dat een agentische aanpak een duidelijke meerwaarde biedt, omdat ze changelogs, release notes en versie-informatie kan samenbrengen tot een bruikbare samenvatting voor ontwikkelaars.

Uit het onderzoek blijkt dat het kernprobleem bij dependencybeheer niet zozeer de detectie van updates is, maar vooral de interpretatie ervan. Door automatische detectie te combineren met een AI-gestuurde analyse verlaagt DepGuard AI de drempel om updates te beoordelen en helpt het systeem ontwikkelaars sneller en beter geïnformeerde beslissingen te nemen. De implementatie en evaluatie tonen daarnaast aan dat de kwaliteit van de uitkomst sterk afhangt van de gekozen orchestratie en de manier waarop broninformatie wordt verwerkt.

Binnen de uitgewerkte oplossingen kwam Mastra als meest geschikte basis naar voren voor verdere uitwerking in de context van Springbok. De combinatie van inhoudelijke kwaliteit, betrouwbaarheid en voldoende aanvaardbare verwerkingstijd maakt deze aanpak het meest interessant voor praktische toepassing. Andere varianten bleken bruikbaar, maar boden minder consistente of minder volledige resultaten.

Deze conclusies moeten wel worden gelezen binnen de grenzen van de uitgevoerde evaluatie. De benchmark was beperkt in omvang en gebaseerd op een beperkt aantal representatieve testgevallen, waardoor verdere validatie nuttig blijft. Toekomstig onderzoek kan zich richten op een bredere evaluatieset, meerdere beoordelingsmodellen en een verdere integratie in een productieomgeving.

De bredere bijdrage van dit werk ligt in de demonstratie dat een agentgebaseerde aanpak meerwaarde kan bieden voor software-onderhoud en technische schuld. In plaats van beslissingen volledig over te nemen, ondersteunt het systeem ontwikkelaars door de interpretatielast te verkleinen en relevante informatie beter toegankelijk te maken.

\clearpage

% \lipsum[76-80]

